% TOP_PROBLEM: {"id": 3323, "title": "Crossing numbers and coloring", "category_id": 3, "category": {"id": 3, "name": "graph_theory", "display_name": "Graph Theory", "description": "Problems involving graphs, networks, and their properties.", "slug": "graph-theory", "order_index": 3, "created_at": "2026-07-31T15:26:25.671Z"}, "status": "open", "problem_number": "OPG-37068", "set_id": 12}
% FIRST_POSTED: 2026-10-01
% ENTRY_KIND: statement_only
% UnsolvedMath ID 3323; source catalogue code OPG-37068
% !TeX program = lualatex
% Definition and sources only; this file is not an LLM solution attempt.
\documentclass[11pt,a4paper]{article}
\usepackage[margin=25mm]{geometry}
\usepackage{fontspec}
\setmainfont{lmroman10-regular.otf}[BoldFont=lmroman10-bold.otf,
 ItalicFont=lmroman10-italic.otf,BoldItalicFont=lmroman10-bolditalic.otf]
\usepackage{amsmath,amssymb,mathtools}
\usepackage{unicode-math}
\setmathfont{latinmodern-math.otf}
\AtBeginDocument{\let\setminus\smallsetminus}
\usepackage{xurl,hyperref}
\hypersetup{colorlinks=true,urlcolor=blue}
\setcounter{secnumdepth}{0}
\setlength{\parindent}{0pt}
\setlength{\parskip}{5pt}
\setlength{\emergencystretch}{3em}
\begin{document}
\section*{Crossing numbers and coloring}\label{3323}
{\small UnsolvedMath ID 3323 · OPG-37068 · Graph Theory · OpenGarden}
\subsection{Definitions and mathematical statement}
Let $G$ be a finite simple undirected graph. Its chromatic number
$\chi(G)$ is the least number of colours in a vertex colouring
where adjacent vertices have different colours. A plane drawing
represents vertices by distinct points and edges by simple arcs
between their ends, with no edge through another vertex, finitely
many transverse crossings of edge interiors, and no triple crossing.
The crossing number $\operatorname{cr}(G)$ is the least total number
of crossing points over all such drawings; endpoints are not counted.
Edges need not be straight.

For an integer $t\ge1$, write $K_t$ for the complete graph on $t$
vertices. Albertson's conjecture is the universal implication
\[
                          \chi(G)\ge t
                 \quad\Longrightarrow\quad
                          \operatorname{cr}(G)\ge\operatorname{cr}(K_t).
\]
The right-hand side is the actual crossing number of $K_t$, not a
conjectured formula for that number. Thus defining this question
does not require solving the separate complete-graph crossing problem.

Equivalently, among graphs with chromatic number exactly $t$, the
complete graph should minimize crossing number. To pass to the
``at least'' formulation, delete vertices until a subgraph of
chromatic number $t$ remains, and use monotonicity of crossing number
under taking subgraphs. The conjecture does not assert that $G$
contains a complete subgraph, a complete subdivision, or a complete
minor. It compares two numerical invariants and permits arbitrarily
many vertices and arbitrary graph structure beyond the chromatic
hypothesis.
\subsection{Short English statement}
Does every graph requiring at least t colours have crossing number at least that of the complete graph on t vertices?
\subsection{Sources}
\begin{itemize}
\item[S1] ULAM.AI and the UnsolvedMath Contributors, supplied problems.json, record 3323 (OPG-37068), OpenGarden collection. Catalogue identity and source transcription; CC BY 4.0.
\url{https://huggingface.co/datasets/ulamai/UnsolvedMath}.
\item[S2] Open Problem Garden, Crossing numbers and coloring. Original problem and discussion; the mathematical formulation below is expanded from the supplied source record.
\url{http://www.openproblemgarden.org/op/crossing_numbers_and_coloring}.
\end{itemize}
\end{document}
